\section{Step 4 Return before accessing an empty matrix}\label{transpose-s04}

Begin the ordinary function by calculating the result count. The dimensions are small nonnegative integers, so their product fits in a signed word. Test both dimensions before the first matrix load. The return label is already supplied below the edit region, so even an unfinished function has a valid return path.

\subsection{Write or inspect}
Append this fragment in the \texttt{naive} region. Reassemble before running. Try zero rows and then zero columns, restoring the default dimensions afterward.

\Needspace{8\baselineskip}
\MARSCodeLabel{MIPS assembly}
\begin{minted}[samepage=true]{mars-idea}
    mul $v0, $a2, $a3
    beq $a2, $zero, naive_return
    nop
    beq $a3, $zero, naive_return
    nop
\end{minted}

\subsection{Expected observation}
Every empty case returns zero without a matrix access. For three by five, the incomplete function returns fifteen but has not copied anything. A correct count alone is insufficient evidence of a transpose.

\Needspace{11\baselineskip}
\subsection{Check your understanding}
Why test the dimensions before the body?

\begin{itemize}
  \item \textbf{A} To avoid the first access when the active extent is empty.
  \item \textbf{B} To force the matrix into the cache.
  \item \textbf{C} To turn every matrix into a square.
\end{itemize}
